\subsection{Baseline and Training Length}\label{sec:baseline}
The generated returns of the ForGAN baseline are reported for the first two configurations of Table~\ref{tab:configurations}, which differ only in their training schedule. They are compared to the statistical properties of the test split from Sec.~\ref{sec:properties-ttf}. Table~\ref{tab:baseline_training_length} reports both configurations alongside the test split. The reference values were computed from the rows of the test window matrix. With $l = 10$, the 401 returns of the test split give $n_{\mathrm{cond}} = 391$ rows. Mode collapse (Sec.~\ref{sec:evaluation}) was not detected in any run of the experimental grid.

\begin{table}[htbp]
  \centering
  \begin{tabular*}{\textwidth}{@{\extracolsep{\fill}}lrrr@{}}
    \toprule
     & Test & Configuration 1 & Configuration 2 \\
    \midrule
    Mean               & $0.0006$  & \ms{-0.0003}{0.0019} & \ms{-0.0029}{0.0019} \\
    Standard deviation & $0.0404$  & \ms{0.0303}{0.0118}  & \ms{0.0365}{0.0034}  \\
    Skewness           & $0.92$    & \ms{-0.18}{0.39}     & \ms{0.19}{0.58}      \\
    Excess kurtosis    & $13.27$   & \ms{0.32}{0.42}      & \ms{4.51}{1.52}      \\
    Minimum            & $-0.1749$ & \ms{-0.1533}{0.0474} & \ms{-0.2234}{0.0287} \\
    Maximum            & $0.3229$  & \ms{0.1357}{0.0519}  & \ms{0.2516}{0.0761}  \\
    $1\%$ quantile     & $-0.1305$ & \ms{-0.0739}{0.0251} & \ms{-0.1074}{0.0162} \\
    $99\%$ quantile    & $0.0903$  & \ms{0.0684}{0.0272}  & \ms{0.1103}{0.0163}  \\
    $P(|x| > 0.05)$    & $0.130$   & \ms{0.115}{0.124}    & \ms{0.129}{0.021}    \\
    $P(|x| > 0.10)$    & $0.026$   & \ms{0.008}{0.016}    & \ms{0.026}{0.008}    \\
    \midrule
    $C_2(1)$           & $0.289$   & \ms{-0.015}{0.056}   & \ms{0.125}{0.087}    \\
    $C_2(2)$           & $0.061$   & \ms{-0.009}{0.036}   & \ms{0.129}{0.055}    \\
    $C_2(5)$           & $0.094$   & \ms{-0.009}{0.035}   & \ms{0.164}{0.107}    \\
    $C_2(10)$          & $0.012$   & \ms{0.018}{0.032}    & \ms{0.073}{0.056}    \\
    \bottomrule
  \end{tabular*}
  \caption{Distributional statistics and autocorrelation of generated
  daily TTF log returns on the test split, ForGAN baseline, configurations 1 and 2}
  \label{tab:baseline_training_length}
\end{table}

\textbf{Heavy tails.} The excess kurtosis of the ForGAN baseline under configuration 1 is \ms{0.32}{0.42}, compared to $13.27$ for the test split. Under configuration 2, the excess kurtosis increases to \ms{4.51}{1.52}, roughly a third of the test split value. The seed spread is large; however, the ranges of the two configurations do not overlap.
The exceedance fractions of configuration 1 do not carry a reliable signal, as the seed spread is larger than the mean, \ms{0.115}{0.124} for $P(|x| > 0.05)$. Under configuration 2 both exceedance fractions match the values of the test split, \ms{0.129}{0.021} against $0.130$ and \ms{0.026}{0.008} against $0.026$. The ForGAN baseline under configuration 2 generates large returns at the observed frequency, but does not generate the few extreme returns that increase the excess kurtosis. 
The standard deviation of the ForGAN baseline under configuration 1 is \ms{0.0303}{0.0118}, compared to $0.0404$ for the test split. Under configuration 2, the standard deviation is \ms{0.0365}{0.0034}, slightly below the test split value.
Longer training increases the tail weight, but the excess kurtosis remains far below the test split. 

\textbf{Asymmetry.} The skewness of the ForGAN baseline under configuration 1 is \ms{-0.18}{0.39}, negative where the test split is positive at $0.92$. Under configuration 2, the skewness is \ms{0.19}{0.58}, roughly a fifth of the test split value. For both configurations, the seed spread crosses zero, and the ranges of the two configurations overlap. As a result, the two configurations cannot be separated based on skewness, unlike the excess kurtosis.
The minimum and maximum values provide additional information about the asymmetry. Under configuration 2 the minimum is \ms{-0.2234}{0.0287} and the maximum is \ms{0.2516}{0.0761}, against $-0.1749$ and $0.3229$ of the test split. The generated distribution extends further downward and less far upward than the test split. Neither configuration of the ForGAN baseline reproduces the skewness of the test split. 

\textbf{Volatility clustering.} The autocorrelation of the ForGAN baseline under configuration 1 is \ms{-0.015}{0.056} for $\tau=1$, compared to $0.289$ for the test split. Therefore, the generated series shows no volatility clustering. Under configuration 2 the autocorrelation is \ms{0.125}{0.087} for $\tau=1$, roughly half of the test split value, indicating some volatility clustering. At $\tau=2, 5, 10$ the autocorrelation under configuration 2 exceeds the test split, resulting in a flat autocorrelation. The generated series does not reproduce the autocorrelation decay described in Sec.~\ref{sec:stylized-facts}. The test split provides a limited reference for the shape of the decay, since its autocorrelation does not decay smoothly either (Sec.~\ref{sec:properties-ttf}).

Table~\ref{tab:mse_training_length} reports the same statistics for the MSE branch, the supervised counterpart of the ForGAN baseline, under configurations 1 and 2. The excess kurtosis of the MSE branch under configuration 1 is \ms{0.19}{0.23}, against \ms{0.32}{0.42} for the ForGAN baseline. Under configuration 2, the excess kurtosis of the MSE branch is \ms{4.13}{1.20}, close to the \ms{4.51}{1.52} of the ForGAN baseline. The two branches show similar patterns in the distributional properties.
The main difference between these two branches is the skewness. The skewness of the MSE branch under configuration 2 is \ms{0.39}{0.27}, compared to \ms{0.19}{0.58} for the ForGAN baseline. The seed spread for MSE does not cross zero, so the MSE branch reproduces a positive skew, unlike the ForGAN baseline. However, the skewness under both configurations is still below the test split value of $0.92$. The MSE branch also varies less across seeds than the ForGAN baseline on the distributional statistics, under both configurations.

\begin{table}[htbp]
  \centering
  \begin{tabular*}{\textwidth}{@{\extracolsep{\fill}}lrrr@{}}
    \toprule
     & Test & Configuration 1 & Configuration 2 \\
    \midrule
    Mean               & $0.0006$  & \ms{0.0025}{0.0026}  & \ms{-0.0011}{0.0020} \\
    Standard deviation & $0.0404$  & \ms{0.0294}{0.0095}  & \ms{0.0354}{0.0011}  \\
    Skewness           & $0.92$    & \ms{-0.07}{0.24}     & \ms{0.39}{0.27}      \\
    Excess kurtosis    & $13.27$   & \ms{0.19}{0.23}      & \ms{4.13}{1.20}      \\
    Minimum            & $-0.1749$ & \ms{-0.1349}{0.0300} & \ms{-0.2151}{0.0396} \\
    Maximum            & $0.3229$  & \ms{0.1387}{0.0513}  & \ms{0.2592}{0.0469}  \\
    $1\%$ quantile     & $-0.1305$ & \ms{-0.0674}{0.0175} & \ms{-0.0976}{0.0083} \\
    $99\%$ quantile    & $0.0903$  & \ms{0.0709}{0.0248}  & \ms{0.1114}{0.0066}  \\
    $P(|x| > 0.05)$    & $0.130$   & \ms{0.103}{0.102}    & \ms{0.124}{0.006}    \\
    $P(|x| > 0.10)$    & $0.026$   & \ms{0.005}{0.010}    & \ms{0.024}{0.002}    \\
    \midrule
    $C_2(1)$           & $0.289$   & \ms{0.040}{0.074}    & \ms{0.149}{0.063}    \\
    $C_2(2)$           & $0.061$   & \ms{0.009}{0.056}    & \ms{0.141}{0.093}    \\
    $C_2(5)$           & $0.094$   & \ms{0.007}{0.055}    & \ms{0.152}{0.088}    \\
    $C_2(10)$          & $0.012$   & \ms{0.024}{0.041}    & \ms{0.134}{0.078}    \\
    \bottomrule
  \end{tabular*}
  \caption{Distributional statistics and autocorrelation of generated
  daily TTF log returns on the test split, MSE branch, configurations 1 and 2}
  \label{tab:mse_training_length}
\end{table}

Configuration 1, corresponding to Algorithm 1 of \textcite{vuleticFinGANForecastingClassifying2024}, leaves the model undertrained on the TTF futures dataset. The generated distribution is close to a normal distribution and shows no volatility clustering.
Configuration 2, corresponding to the published implementation of \textcite{vuleticFinGANForecastingClassifying2024}, is a usable base configuration for training on the TTF futures dataset. However, the excess kurtosis and skewness remain well below the test split values.
The two branches respond to the schedule in the same way, so the difference between the two configurations is a result of the training length and not the objective.